% id: 108-36
% book: 쎈B 대수 · 08 등차수열과 등비수열
% section: 유형 08 두 수 사이에 수를 넣어서 만든 등차수열의 합
% figure: none
% answer: 
% answer_source: 
% variant_level: 0 (원본 전사)
% 이 파일은 build.mjs 가 items.json 에서 만든다. 고칠 때는 items.json 을 고친다.
\smash{\raisebox{1.5mm}{\dmkichul{쎈B 대수 108쪽 36번}}}
수열 $15$, $a_1$, $a_2$, $a_3$, $\cdots$, $a_n$, $-19$가 이 순서대로 등차수열을 이루고 $a_1+a_2+a_3+\cdots+a_n=-32$일 때, 자연수 $n$의 값은?
\dmchoicesauto{}{$15$}{$16$}{$17$}{$18$}{$19$}
